% =====================================================================
% Method: Linear Representational Alignment
% =====================================================================
\section{Method: Linear Representational Alignment}
\label{sec:method}

The transplant of Eq.~\eqref{eq:transplant} needs, for every adapted module
$\ell$, two matrices $P^{\mathrm{in}}_\ell$ and $P^{\mathrm{out}}_\ell$ that
translate between the activation geometries of $M_v$ and $M_{v+1}$. This section
specifies how we estimate them from anchors alone, why we estimate them in
activation space rather than weight space, and the variants we ablate.

\subsection{Anchor activations}
\label{sec:method:anchors}

We draw $\mathcal{D}_{\mathrm{anchor}}$ from a generic corpus that is unrelated
to any downstream task (in our experiments, WikiText-2 training text), tokenize
it once with the shared tokenizer, and run the same token ids through both base
checkpoints with forward hooks on every module in $\mathcal{L}$. For module
$\ell$ this yields token-aligned matrices
\[
  H^{\mathrm{in},v}_\ell,\;H^{\mathrm{in},v+1}_\ell\in\mathbb{R}^{N\times d^{\mathrm{in}}_\ell},
  \qquad
  H^{\mathrm{out},v}_\ell,\;H^{\mathrm{out},v+1}_\ell\in\mathbb{R}^{N\times d^{\mathrm{out}}_\ell},
\]
where row $i$ of each matrix is the activation of the same token position in
the same sequence. Pad positions are dropped and, to bound memory, the kept
positions are subsampled with a fixed seed so both models keep identical rows.
Because the inputs are identical, the $i$-th rows are two views of one
underlying token in two coordinate systems, which is what makes a regression
between them meaningful. This is the only pass through the models the method
requires; the activations are cached and reused by every variant below.

\subsection{Closed-form maps}
\label{sec:method:maps}

Each map is a least-squares regression from one version's activations to the
other's. Writing $S$ for the source matrix and $T$ for the target (both
mean-centered over rows), we solve for $P$ with $S P^\top\approx T$:
\begin{align}
  \text{ridge:}\quad
    &P=\arg\min_P\;\lVert T-SP^\top\rVert_F^2+\lambda\lVert P\rVert_F^2
     \;=\;(T^\top S)(S^\top S+\lambda I)^{-1},
  \label{eq:ridge}\\
  \text{orthogonal (Procrustes):}\quad
    &P=\arg\min_{P^\top P=I}\;\lVert T-SP^\top\rVert_F^2
     \;=\;(UV^\top)^\top,\quad S^\top T=U\Sigma V^\top .
  \label{eq:procrustes_sol}
\end{align}
The input map takes $S=H^{\mathrm{in},v+1}_\ell$, $T=H^{\mathrm{in},v}_\ell$
(it pulls a new-version input back into the old basis the adapter was trained
to read), and the output map takes $S=H^{\mathrm{out},v}_\ell$,
$T=H^{\mathrm{out},v+1}_\ell$ (it pushes the adapter's old-basis output forward
into the new basis). Both are $d\times d$ solves per module, dominated by one
$d^3$ inverse or SVD; for $d\le 4096$ the complete set of maps for a model is
fit in seconds on a CPU. Ridge is the default. The orthogonal variant encodes
the prior that version drift is a rotation; \S\ref{sec:ablations} shows that
prior is wrong often enough to cost accuracy, which is a central empirical
point of the paper.

\subsection{Transplant and loading}
\label{sec:method:transplant}

With the maps in hand the adapter is rewritten factor by factor,
\[
  A'_\ell = A_\ell P^{\mathrm{in}}_\ell,\qquad
  B'_\ell = P^{\mathrm{out}}_\ell B_\ell,
\]
so that $\tfrac{\alpha}{r}B'_\ell A'_\ell=P^{\mathrm{out}}_\ell\Delta W_\ell
P^{\mathrm{in}}_\ell$. The rank, scale $\alpha$, dropout, and target-module
list are unchanged, so the result is a valid adapter in the original
serialization format and loads onto $M_{v+1}$ with standard PEFT tooling. No
new parameters are introduced and inference cost is identical to the original
adapter. Modules that were not adapted need no map.

\paragraph{One map serves every adapter.}
The maps depend only on the two base checkpoints, not on the adapter, so they
can be fit once per version pair and applied to an entire adapter library. This
costs nothing in fit quality on the input side. Suppose instead we fit the
adapter's read-out directly, choosing $A'$ to minimize
$\lVert H^{\mathrm{in},v+1}A'^\top-H^{\mathrm{in},v}A^\top\rVert_F^2+\lambda\lVert A'\rVert_F^2$,
so that the transplanted adapter computes the same $r$ features on the new
model as the original did on the old. The solution is
$A'^\top=(S^\top S+\lambda I)^{-1}S^\top TA^\top$ with $S,T$ as in
Eq.~\eqref{eq:ridge}, which is exactly $A P^{\mathrm{in}}$ for the ridge map:
least squares decomposes over target coordinates, so the adapter-agnostic map
is already the adapter-optimal one for every adapter and every $\lambda$. (We
confirm this numerically to float precision.) A vendor could therefore publish
$\{P^{\mathrm{in}}_\ell,P^{\mathrm{out}}_\ell\}$ alongside a release as a
compatibility artifact, with no per-adapter refitting.

\begin{algorithm}[t]
\caption{Activation-aligned adapter transplant}
\label{alg:transplant}
\begin{algorithmic}[1]
\Require bases $M_v$, $M_{v+1}$; adapter $\{(A_\ell,B_\ell)\}_{\ell\in\mathcal{L}}$;
         unlabeled anchors $\mathcal{D}_{\mathrm{anchor}}$; solver $\in\{$ridge($\lambda$), orthogonal$\}$
\State tokenize $\mathcal{D}_{\mathrm{anchor}}$ once; run the ids through $M_v$ and $M_{v+1}$ with hooks on $\mathcal{L}$
\State cache token-aligned $H^{\mathrm{in},\cdot}_\ell$, $H^{\mathrm{out},\cdot}_\ell$ for every $\ell\in\mathcal{L}$
\For{$\ell\in\mathcal{L}$}
  \State $P^{\mathrm{in}}_\ell\gets\textsc{Fit}(S{=}H^{\mathrm{in},v+1}_\ell,\;T{=}H^{\mathrm{in},v}_\ell)$ \Comment{Eq.~\eqref{eq:ridge} or \eqref{eq:procrustes_sol}}
  \State $P^{\mathrm{out}}_\ell\gets\textsc{Fit}(S{=}H^{\mathrm{out},v}_\ell,\;T{=}H^{\mathrm{out},v+1}_\ell)$
  \State $A'_\ell\gets A_\ell P^{\mathrm{in}}_\ell$;\quad $B'_\ell\gets P^{\mathrm{out}}_\ell B_\ell$
\EndFor
\State \Return adapter $\{(A'_\ell,B'_\ell)\}$ loaded onto $M_{v+1}$
\end{algorithmic}
\end{algorithm}

\subsection{Why activations rather than weights}
\label{sec:method:why}

The alternative, used by LoRA-X and Cross-LoRA, is to relate the two versions
through their weight matrices. Three properties of version drift argue for
the activation route.
First, the adapter acts on activations: $A_\ell$ reads coordinates of
$x\in\mathbb{R}^{d^{\mathrm{in}}}$ and $B_\ell$ writes coordinates of the module
output. The quantity that must be preserved across versions is the function
$x\mapsto\Delta W_\ell x$ \emph{on the inputs the module actually receives},
and the empirical distribution of those inputs is available only by running
the model. A regression on anchor activations minimizes exactly the error in
that function; a weight-subspace projection minimizes a surrogate.
Second, a module's dominant singular directions are a property of its weight
alone, whereas the drift we need to undo is a joint property of the whole
network upstream of the module. Two versions can share nearly identical
dominant subspaces for $W_\ell$ while feeding $W_\ell$ a differently oriented
input distribution; weight-space methods are blind to this, activation methods
see it directly (\S\ref{sec:drift} quantifies how much of the drift is of this
kind).
Third, drift need not be orthogonal, nor confined to a low-rank subspace.
Ridge regression returns a general linear map; an orthogonal map or a
subspace projection cannot represent anisotropic rescaling, and we show in
\S\ref{sec:ablations} that this restriction alone forfeits a measurable
fraction of the recoverable gap.

\subsection{The Cross-LoRA baseline in the equal-dimension case}
\label{sec:method:crosslora}

Cross-LoRA~\citep{xia2025crosslora} computes rank-$k$ truncated SVDs
$W_s\approx U_s\Sigma_sV_s^\top$ and $W_t\approx U_t\Sigma_tV_t^\top$ of the
source and target base weights of each module, fits
$\hat P_U=\arg\min_P\lVert PU_s-U_t\rVert_F$ and
$\hat P_V=\arg\min_P\lVert PV_s-V_t\rVert_F$, sets $\tilde U_s=\hat P_UU_s$,
$\tilde V_s=\hat P_VV_s$, and transfers
$\Delta W_t=\tilde U_s(\tilde U_s^\top\Delta W_s\tilde V_s)\tilde V_s^\top$.
When $W_s$ and $W_t$ have the same shape, which is the case for every version
pair we study, $U_s$ has orthonormal columns and the minimum-norm least-squares
solution is $\hat P_U=U_tU_s^\top$, so $\tilde U_s=U_t$ exactly and likewise
$\tilde V_s=V_t$. The transfer therefore reduces to
\begin{equation}
  \Delta W_t \;=\; (U_tU_t^\top)\,\Delta W_s\,(V_tV_t^\top),
  \label{eq:crosslora_equal}
\end{equation}
the projection of the source update onto the top-$k$ left and right singular
subspaces of the \emph{target} weight. As $k\to\min(d^{\mathrm{in}},d^{\mathrm{out}})$
this approaches naive transfer. We implement the method literally (truncated
SVD, least-squares alignment, projection) and verify numerically that it
matches Eq.~\eqref{eq:crosslora_equal}; like ours it preserves rank
($B'=U_tU_t^\top B$, $A'=AV_tV_t^\top$) and adds no inference cost, so the
comparison isolates the alignment signal, weights versus activations, with
everything else held fixed. We use the paper's truncation rank $k=320$ and
sweep it in \S\ref{sec:ablations}.

\subsection{Variants}
\label{sec:method:variants}

We ablate four design choices. \textbf{Solver}: ridge with $\lambda$ swept
over four orders of magnitude, versus the orthogonal solution.
\textbf{Scope}: one pair of maps per module (default) versus a single pair
shared by all modules of equal shape, which tests whether drift is a global
rotation of the residual stream or module-specific. \textbf{Anchor budget}:
the number of anchor tokens $N$ used to fit the maps, from $10^2$ to $10^4$.
\textbf{Module subset}: aligning only the attention projections with the
largest measured drift versus all adapted modules. A nonlinear (small MLP)
map is a natural extension that would break rank preservation and is left to
future work.
